\chapter{Full list of codes}
\label{app:code}

This appendix collects the complete code of the core modules implemented in Volume 1. Although the text is partially excerpted for explanation purposes, the complete code is provided here. Check the GitHub repository for the latest version.

\section{project structure}

\begin{lstlisting}[style=bashstyle]
make-llm-basic/
  src/makellm/
    __init__.py
    tokenizer/
      __init__.py
      base.py        # Tokenizer abstract class
      bpe.py         # BPE tokenizer
      unigram.py     # Unigram tokenizer
    model/
      __init__.py
      embedding.py   # TokenEmbedding, PositionalEmbedding
      attention.py   # MultiHeadAttention, CausalMask
      transformer_block.py
      gpt.py         # GPT model
    training/
      __init__.py
      dataset.py     # LMDataset
      loss.py        # cross_entropy_loss
      optimizers.py  # AdamW + cosine scheduler
      trainer.py     # Trainer class
    inference/
      __init__.py
      sampler.py     # Greedy, Top-K, Top-P, Temperature
      generate.py    # generate function
    utils/
      __init__.py
      config.py      # ModelConfig, TrainerConfig
      seed.py        # set_seed
      logging.py     # Logger
  tests/             # pytest unit testing
  book/              # of this book LaTeX sauce
\end{lstlisting}

\section{run test}

\begin{lstlisting}[style=bashstyle]
# install package
cd make-llm-basic
pip install -e .

# Run full test
pytest tests/ -v

# Test only specific modules
pytest tests/test_tokenizer.py -v
pytest tests/test_model.py -v
pytest tests/test_training.py -v
\end{lstlisting}

\section{Quickstart example}

\begin{lstlisting}[language=Python]
"""Using all modules from Volume 1 end-to-end example."""
from makellm.tokenizer import BPETokenizer
from makellm.model import GPT
from makellm.training import LMDataset, Trainer
from makellm.inference import generate, TopPSampler
from makellm.utils import ModelConfig, TrainerConfig, set_seed

set_seed(42)

# 1. Tokenizer training
corpus = ["the quick brown fox jumps", "the lazy dog sleeps"] * 100
tokenizer = BPETokenizer()
tokenizer.train(corpus, vocab_size=500)

# 2. Model creation
model_config = ModelConfig(
    vocab_size=tokenizer.vocab_size,
    context_length=32, d_model=64, n_heads=4, n_layers=2, d_ff=128
)
model = GPT(model_config)
print(f"Model: {model}")
print(f"Parameters: {model.num_parameters():,}")

# 3. Dataset preparation
dataset = LMDataset(corpus, tokenizer, context_length=32)

# 4. learning
trainer_config = TrainerConfig(
    batch_size=4, learning_rate=3e-4, max_steps=100,
    warmup_steps=10, device="cpu"
)
trainer = Trainer(model, dataset, trainer_config, model_config)
trainer.train()

# 5. Text creation
sampler = TopPSampler(p=0.9, temperature=0.8)
text = generate(model, tokenizer, prompt="the", max_new_tokens=20, sampler=sampler)
print(f"Generated: {text}")
\end{lstlisting}

\section{By module API summation}

\begin{center}
\begin{tabularx}{\textwidth}{lX}
\toprule
\textbf{module}&\textbf{main API}\\
\midrule
\code{makellm.tokenizer} & \code{BPETokenizer, UnigramTokenizer} \\
\code{makellm.model} & \code{GPT, MultiHeadAttention, TransformerBlock} \\
\code{makellm.training} & \code{LMDataset, Trainer, cross\_entropy\_loss} \\
\code{makellm.inference} & \code{generate, GreedySampler, TopKSampler, TopPSampler} \\
\code{makellm.utils} & \code{ModelConfig, TrainerConfig, set\_seed} \\
\bottomrule
\end{tabularx}
\end{center}

\section{Setting data class}

\begin{lstlisting}[language=Python]
@dataclass
class ModelConfig:
    vocab_size: int = 1000
    context_length: int = 128
    d_model: int = 128
    n_heads: int = 4
    n_layers: int = 4
    d_ff: int = 512
    dropout: float = 0.1
    use_bias: bool = True
    layer_norm_eps: float = 1e-5
    pad_token_id: int = 0

@dataclass
class TrainerConfig:
    batch_size: int = 32
    learning_rate: float = 3e-4
    weight_decay: float = 0.1
    max_epochs: int = 10
    warmup_steps: int = 100
    lr_scheduler: str = "cosine"
    grad_clip: float = 1.0
    log_every: int = 10
    eval_every: int = 200
    save_every: int = 1000
    seed: int = 42
    device: str = "cpu"
    out_dir: str = "./runs"
\end{lstlisting}

By adjusting these settings, you can experiment with models of different sizes. For example, a similar setup to GPT-2 small is\code{d\_model=768, n\_heads=12, n\_layers=12, context\_length=1024}(but actual training requires significant computing resources).

\section{license}

The code is released under the MIT license. It can be freely copied, modified, and used commercially. The text of the book is CC BY-NC-SA 4.0 (author verification required).
